\chapter*{Introduction Générale}
\addcontentsline{toc}{chapter}{Introduction Générale}

L'évolution fulgurante des technologies de l'information et de la communication a radicalement transformé le paysage éducatif mondial \cite{ai_education}. Au cours de la dernière décennie, nous avons assisté à une transition massive des salles de classe physiques vers des environnements d'apprentissage hybrides et entièrement virtuels. Cependant, cette transition ne s'est pas faite sans heurts : les institutions académiques sont aujourd'hui confrontées à des défis majeurs tels que l'engagement des étudiants à distance, la personnalisation de l'apprentissage et la surcharge administrative des enseignants.

\section*{Le Contexte de la Transformation Digitale}
La pandémie de COVID-19 a accéléré l'adoption du e-learning, mais a également révélé les limites des outils existants. Les enseignants se retrouvent souvent isolés face à des outils complexes, tandis que les étudiants perdent leur motivation devant des interfaces peu interactives. La nécessité d'une plateforme "intelligente" capable d'accompagner l'étudiant dans sa révision tout en simplifiant la tâche du professeur est devenue une priorité pour les universités d'élite comme l'EPI.

\section*{Objectifs et Problématique}
Les systèmes de gestion de l'apprentissage (LMS) traditionnels se limitent souvent à être de simples dépôts de fichiers statiques. Parallèlement, les outils de collaboration génériques comme Microsoft Teams, bien que puissants, manquent d'une couche pédagogique spécifique intégrée. C'est dans ce contexte que naît le projet \textbf{eduGenius}.
...

\textbf{eduGenius} se positionne comme une plateforme "SaaS" d'éducation nouvelle génération, fusionnant les capacités collaboratives d'un outil de communication professionnel avec la puissance analytique et générative de l'Intelligence Artificielle. Conçue pour \textbf{Polytech-Intl (École Polytechnique Internationale)}, cette plateforme a été développée au sein de l'organisme d'accueil \textbf{BES2I (Business Expert Software \& Systems Integrator)}. Elle vise à offrir un écosystème complet où la visioconférence \cite{daily_co}, le chat en temps réel \cite{socket_io}, la gestion de contenu et l'évaluation automatisée cohabitent de manière fluide.

L'innovation majeure d'eduGenius réside dans son exploitation des \textbf{Large Language Models (LLM)} comme Google Gemini \cite{google_gemini}. Pour l'étudiant, l'IA devient un tuteur personnel capable de synthétiser des cours complexes et de générer des outils de mémorisation. Pour l'enseignant, elle agit comme un assistant pédagogique réduisant drastiquement le temps de conception des évaluations.

Ce rapport retrace le cycle de vie complet de ce projet, de la phase d'incubation conceptuelle à la mise en production technique, en suivant la méthodologie Scrum \cite{scrum_guide}. Le manuscrit est structuré comme suit :

\begin{itemize}
    \item \textbf{Chapitre 1 : Cadre Général du Projet} - Ce chapitre définit le périmètre d'intervention, présente l'EPI en tant qu'organisme d'accueil, analyse l'existant et expose la problématique centrale.
    \item \textbf{Chapitre 2 : État de l'art} - Une étude approfondie des fondements théoriques des technologies modernes du Web (MERN Stack \cite{mern_stack}, Next.js 16 \cite{nextjs_docs}) et des concepts d'IA Générative appliqués à l'éducation.
    \item \textbf{Chapitre 3 : Analyse et Spécification des Besoins} - Une phase de modélisation rigoureuse identifiant les acteurs et capturant les besoins fonctionnels et techniques.
    \item \textbf{Chapitre 4 : Conception} - Ce chapitre détaille l'architecture logicielle, la conception de la base de données NoSQL \cite{mongodb_manual} et les interactions dynamiques.
    \item \textbf{Chapitre 5 : Réalisation} - Une présentation de l'environnement de développement, de l'implémentation technique des modules clés et une démonstration des interfaces.
\end{itemize}

Enfin, une conclusion générale synthétisera les résultats obtenus et ouvrira des perspectives sur l'avenir de l'IA dans l'éducation.
